\chapter*{Appendix C: Individual Reflection --- Ali Hamza}
\addcontentsline{toc}{chapter}{Appendix C: Individual Reflection --- Ali Hamza}
\textbf{Name:} Ali Hamza \hfill \textbf{Banner ID:} [B00XXXXXX] \\[4pt]
\textbf{Role:} Backend Developer --- Node.js, Express.js \& MongoDB API
\vspace{1em}

\section*{C.1 Self-Reflection}

Before this project, my backend experience was limited to simple CRUD applications with synchronous request-response cycles. The architecture demanded here --- asynchronous job queuing, JWT-based stateless authentication with refresh token rotation, GDPR-compliant cascading deletes, and integration with two external Python microservices --- was a significant step up in complexity and professional relevance.

The most transformative learning was in asynchronous system design. I initially implemented NLP processing as synchronous HTTP calls from the resume upload endpoint: the controller called the NLP parser, waited for the response, then called the matcher, and only then returned to the client. This worked in isolation but failed under any realistic load. After reading about the Bull queue pattern and implementing exponential backoff retry logic, I understood why asynchronous job queuing is standard practice in production systems. The insight that a web server's job is to accept work and return quickly --- not to do the work itself --- changed how I think about API design.

JWT implementation revealed subtleties I had not anticipated. Generating access tokens is straightforward; managing refresh token rotation without creating race conditions under concurrent requests is not. When two requests arrive within the same access-token expiry window, both attempt to refresh simultaneously, and without careful handling, the second refresh invalidates the token generated by the first. Solving this with the HttpOnly cookie approach and a short-lived access token expiry of 15 minutes required reading RFC 7519 carefully and testing edge cases deliberately rather than assuming my implementation was correct.

GDPR compliance was a dimension of backend development I had not previously considered. Implementing the \texttt{deleteAccount} endpoint --- which must cascade delete across Resume, Match, and AuditLog collections before removing the User document --- required thinking through the failure modes: what if a network error interrupts the cascade mid-way? I implemented a compensating transaction pattern using ordered delete operations, but acknowledged in comments that a proper distributed saga pattern would be needed for production reliability.

Security hardening with Helmet.js and rate limiting was intellectually satisfying. Configuring the fourteen HTTP security headers set by Helmet --- understanding what each one does and why it matters --- gave me a concrete, applied understanding of web security that lectures had conveyed only abstractly. The experience of reading the OWASP Top Ten \citep{OWASP2021} and then directly mapping each risk to a mitigation in my codebase (express-validator for injection, requireRole middleware for broken access control, Helmet for misconfiguration) was genuinely educational.

Where I fell short was in testing coverage. I wrote unit tests for the auth controller and queue service, but integration tests covering the full upload-parse-match pipeline were not completed within the project timeline. This means the error-handling paths --- particularly the Bull job failure state transition to \texttt{'failed'} after three retries --- were validated manually rather than automatically. This represents a meaningful gap in quality assurance.

\textbf{Key strengths demonstrated:} end-to-end async pipeline design; JWT auth with refresh token rotation; GDPR-aware data deletion cascade; OWASP security hardening.

\textbf{Skills developed:} Bull/Redis job queue configuration; Mongoose schema design with indexes; express-validator middleware chains; Morgan HTTP logging; Docker multi-stage builds for Node.js.

\section*{C.2 Critical Appraisal}

The backend API delivers all specified endpoints across four route groups (auth, jobs, resumes, matches) with JWT role-based access control, rate limiting, input validation, and structured error responses. The Bull queue integration with exponential backoff and explicit failure state transitions meets the supervisor's requirement for fault-tolerant asynchronous processing.

A significant architectural limitation is the absence of a message broker health check in the startup sequence. If Redis is unavailable when the Node.js server starts, Bull silently fails to initialise and resume uploads are accepted but never processed. A startup health check with a fast-fail and clear error message would prevent this silent failure mode in deployment.

MongoDB schema design choices also merit scrutiny. The Resume model stores the full \texttt{rawText} field (potentially tens of kilobytes per document) alongside structured \texttt{parsedProfile} data. For a corpus of tens of thousands of resumes, this creates substantial storage overhead. Separating raw text into a GridFS store or a dedicated collection with a reference from the Resume document would be the appropriate pattern at scale.

The role-based access control implementation is binary (recruiter or candidate) with no support for administrative users, organisation-level multi-tenancy, or fine-grained permissions. Real recruitment platforms require at minimum a third role (organisation admin) and the ability to scope jobs and candidates to organisational accounts. Extending the RBAC model to support this would be the primary architecture change needed to move towards commercial viability.

Rate limiting at 200 requests per 15 minutes per IP is a reasonable research-prototype setting, but IP-based rate limiting is trivially bypassed in practice and does not protect against authenticated users abusing the API. Token-bucket rate limiting at the user-id level would be more appropriate for a production system. Despite these limitations, the backend provides a solid, documented foundation for the full-stack integration and demonstrates professional-level API design practices.
